\chapter{Aims, Objectives, and Research Questions}
\label{ch:aims}

This chapter states the purpose of the thesis in its final form. The project began as a Docker-based teaching artefact and evolved into a research investigation of service-time load dependence in containerised microservices. The aims and objectives below reflect both roles.

\section{Overall Aim}

The overall aim is to develop and validate a reproducible capacity-modelling framework for Dockerised request-response services, capable of showing when classical M/G/1 and M/G/$c$ assumptions hold, when they fail, and what modelling correction is needed when measured service time changes with load.

The framework is intended to serve two audiences:

\begin{itemize}
    \item \textbf{Teaching}: COMP9334 students should be able to observe queueing theory, saturation, response-time CDFs, and DES validation using real services on a personal machine.
    \item \textbf{Research and capacity planning}: performance engineers should be able to use request-level traces to diagnose whether a microservice is approximately constant-service, load-dependent, structurally mismatched, or affected by runtime artefacts.
\end{itemize}

\section{Research Questions}

The thesis is organised around five research questions.

\begin{description}
    \item[RQ1.] Can a lightweight Docker platform generate reproducible request-level traces suitable for queueing and DES validation?
    \item[RQ2.] When do standard M/G/1 and M/G/$c$ models accurately predict measured response-time distributions for containerised services?
    \item[RQ3.] Do CPU-bound containerised services preserve the constant-service-time assumption as arrival rate and utilisation increase?
    \item[RQ4.] If service time is load-dependent, does a fitted service-time model improve DES response-time CDFs relative to the constant-service baseline?
    \item[RQ5.] Are richer explicit scheduling models, such as priority/preemptive-resume or feedback queues, supported by the current traces, or do they require new experiments?
\end{description}

\section{Objectives}

To answer these questions, the thesis pursues the following concrete objectives.

\begin{enumerate}
    \item Build a Docker Compose platform containing multiple server implementations across different runtimes, workloads, and concurrency models.
    \item Implement open-loop Poisson load generation and standardise the trace format to include \texttt{arrival\_unix\_ns}, \texttt{service\_ms}, \texttt{queue\_ms}, \texttt{response\_ms}, and \texttt{status\_code}.
    \item Implement trace-driven M/G/1 and M/G/$c$ DES models using replayed arrival times, lognormal service-time sampling, and response-CDF validation.
    \item Fit and compare M0, M1, and M2 service-time models, using AIC and response-CDF KS distance to determine whether service-time modelling improves response-time prediction.
    \item Inspect service-time distributions using histograms and test service-time independence from arrival rate using ANOVA on $\log(\texttt{service\_ms})$.
    \item Identify structural mismatch cases where a single shared M/G/$c$ queue is not the correct abstraction, including Node.js cluster and Go SQLite.
    \item Implement exploratory M3, priority/preemptive-resume, and feedback-priority models to test whether the load-dependent finding can be replaced or strengthened by an explicit scheduling interpretation.
    \item Document the full research process, including failed assumptions and debugging discoveries, so that the final model is presented as an empirical discovery rather than as a post-hoc curve fit.
\end{enumerate}

\section{Scope}

The thesis focuses on single-machine Docker experiments under controlled load. It does not claim production-scale validation, Kubernetes deployment, or cloud-native autoscaling. The results are therefore strongest as empirical evidence for the existence and diagnosis of load-dependent service time in controlled containerised services.

The thesis also distinguishes between completed work and future work. M/G/1, M/G/$c$, M0/M1/M2, M3, ANOVA, and exploratory priority/feedback validation are implemented and evaluated. Fixed-point forecasting, arrival-rate prediction, real two-class priority experiments, tandem-queue modelling for SQLite, and cloud/native-Linux validation remain future-work directions beyond the current submission.

\section{Claim Boundaries}

The core claim is not that one model is universally best. The evidence instead supports a diagnostic workflow:

\begin{enumerate}
    \item If replay DES fits well, the queueing structure is likely adequate.
    \item If M0 fails but M2 improves the response CDF, service time is likely load-dependent.
    \item If all one-stage models fail, the server may contain a structural mismatch such as split queues, tandem queues, runtime warmup, or hidden locks.
    \item If explicit priority or feedback models are desired, the experiment must log real priority classes or feedback visits; they should not be inferred from one-class traces alone.
\end{enumerate}

This boundary is important for thesis-level defensibility. The work contributes a tested method for finding model violations, not a universal replacement for queueing theory.
